\documentclass[../../main.tex]{subfiles}
\begin{document}

\begin{flushright}
\footnotesize\textit{【自動文字起こし・要確認】}
\end{flushright}

{\bf [解]}
\begin{align*}
f(x)>0\ (0\le x\le1), \qquad f(0)=2,\ f(1)=1 \quad\cdots\text{①}
\end{align*}
又，題意及び①から$k\in\mathbb{R}_{>0}$として，
\begin{align*}
(x-a)\{f(a)+f(x)\}=k\int_a^xf(t)dt \quad\cdots\text{②}
\end{align*}
が$0\le a\le x\le1$なる任意の$a,x$で成立．（$a=x$でも明らかに成立．）$x$で微分して，
\begin{align*}
f(a)+f(x)+xf'(x)-af'(x)=kf(x)
\end{align*}
$y=f(x)$とおく．$a=x=1$として①から$k=2$．次に$a=0$として
\begin{align*}
2+y+x\dfrac{dy}{dx}=2y
\end{align*}
$y\ne2,\ x\ne0$の時
\begin{align*}
\dfrac{dy}{y-2}=\dfrac{dx}{x} \quad\cdots\text{③}
\end{align*}
積分して，
\begin{align*}
\log|y-2|=\log x+C_1 \quad\cdots\text{④}
\end{align*}
ただし，$C_1$は定数．$x=1$は①より，$x\ne0,\ y\ne2$をみたすので，$x=1$で④が成立することから$C_1=0$．したがって
\begin{align*}
x=|y-2|
\end{align*}
したがって，$(x,y)=(1,1)$とあわせて
\begin{align*}
y=-x+2
\end{align*}
これは$x=0$でも成立する．


\newpage

\end{document}
